\pagestyle{empty}
\tight
\begin{tblr}{width=\textwidth,colspec={|Q[c,m,1.9cm]|X[l,m]|},hlines,vlines,hspan=minimal,row{1}={bg=headblue},rowsep=1pt,colsep=3pt}
\SetCell{c,m}\hfil\logo\hfil & \textbf{POLITEKNIK NEGERI MALANG}\par\textbf{JURUSAN TEKNOLOGI INFORMASI}\par\textbf{PROGRAM STUDI: D4 TEKNIK INFORMATIKA} \\
\SetCell[c=2]{c} \textbf{RENCANA PEMBELAJARAN SEMESTER (RPS)} & \\
\end{tblr}
\begin{longtblr}{width=\textwidth,colspec={|Q[l,m,1.9cm]|X[0.85,l,m]|X[1.45,l,m]|X[0.75,l,m]|X[1.15,l,m]|},hlines,vlines,hspan=minimal,rowsep=1pt,colsep=2pt,row{1,3}={bg=headgray,font=\bfseries}}
MATA KULIAH & KODE & BOBOT (SKS)/jam & SEMESTER & TGL. PENYUSUNAN \\
Kecerdasan Artifisial & RTI254004 & 2 SKS/4 jam & 4 & 1 Juli 2025 \\
OTORISASI & Dosen Pengembang RPS & Koordinator RMK & \SetCell[c=2]{l} Ka PRODI &  \\
 & Vivi Nur Wijayaningrum, S.Kom., M.Kom. &  & \SetCell[c=2]{l}{\vspace{8mm}\par Dr. Ely Setyo Astuti, S.T., M.T.} &  \\
\SetCell[r=10]{l} Capaian Pembelajaran (CP) & \SetCell[c=4]{l,bg=headgray,font=\bfseries} Capaian Pembelajaran Lulusan Program Studi (CPL-Prodi) &  &  &  \\
 & CPL07 & \SetCell[c=3]{l} Mampu mengembangkan sistem komputasi cerdas dan melakukan analisis data berbasis kecerdasan buatan untuk pengambilan keputusan. &  &  \\
 & CPL09 & \SetCell[c=3]{l} Mampu menjelaskan cara kerja sistem komputer dan menerapkan pendekatan komputasi untuk menyelesaikan masalah. &  &  \\
 & CPL10 & \SetCell[c=3]{l} Mampu menganalisis persoalan kompleks dengan wawasan multidisiplin untuk menghasilkan solusi di bidang informatika dan ilmu komputer. &  &  \\
 & \SetCell[c=4]{l,bg=headgray,font=\bfseries} Capaian Pembelajaran mata kuliah (CPL-MK) &  &  &  \\
 & CPMK0706 & \SetCell[c=3]{l} Mahasiswa mampu mengembangkan sistem cerdas berbasis AI, machine learning, dan analisis data untuk menghasilkan keputusan yang bermakna. &  &  \\
 & CPMK1005 & \SetCell[c=3]{l} Mahasiswa mampu mengevaluasi solusi sistem cerdas berdasarkan kinerja, keandalan, dan interpretabilitas. &  &  \\
 & \SetCell[c=4]{l,bg=headgray,font=\bfseries} Kemampuan akhir tiap tahapan belajar (Sub-CPMK) &  &  &  \\
 & SCPMK0706-03101 & \SetCell[c=3]{l} Mahasiswa mampu mengimplementasikan algoritma pencarian, representasi pengetahuan, dan sistem inferensi berbasis aturan. &  &  \\
 & SCPMK1005-03102 & \SetCell[c=3]{l} Mahasiswa mampu mengevaluasi dan membandingkan kinerja berbagai algoritma AI berdasarkan metrik evaluasi yang sesuai dengan konteks permasalahan. &  &  \\
\SetCell[r=2]{l} Penilaian & \SetCell[c=4]{l} *Nilai CPMK didapat dari agregasi nilai SUB-CPMK yang dibebankan &  &  &  \\
 & \SetCell[c=4]{l}{\tiny\begin{tblr}{width=\linewidth,colspec={|Q[c,m,0.1\linewidth]|Q[c,m,0.16\linewidth]|X[l,m]|X[l,m]|X[l,m]|X[l,m]|X[l,m]|X[l,m]|Q[c,m,0.08\linewidth]|},hlines,vlines,hspan=minimal,rowsep=1pt,colsep=0.7pt,row{1,2}={bg=headgray,font=\bfseries}}
Kode CPMK & Kode SCPMK & \SetCell[c=6]{c} Bobot per Bentuk Penilaian & & & & & & Total Bobot \\
 & & Tugas & Kuis 1 & UTS & Kuis 2 & PBL & UAS & \\
CPMK0706 & SCPMK0706-03101 & 18\%: implementasi algoritma AI & 5\%: pencarian dan agen & 10\%: pencarian dan representasi pengetahuan & 0\% & 12\%: implementasi sistem AI & 10\%: presentasi sistem AI & 55\% \\
CPMK1005 & SCPMK1005-03102 & 12\%: evaluasi algoritma AI & 0\% & 5\%: representasi pengetahuan & 5\%: sistem pakar dan ML & 10\%: evaluasi sistem AI & 13\%: evaluasi dan presentasi & 45\% \\
\SetCell[c=8]{r}\textbf{Total Per Penilaian} & & & & & & & & \textbf{100\%} \\
\end{tblr}} &  &  &  \\
Diskripsi Singkat Mata Kuliah & \SetCell[c=4]{l} Kecerdasan Artifisial membangun pemahaman dan kemampuan mengimplementasikan algoritma pencarian, representasi pengetahuan, sistem pakar, dan pengantar ML untuk menyelesaikan permasalahan komputasi cerdas. &  &  &  \\
Bahan Kajian / Materi Pembelajaran & \SetCell[c=4]{l}\begin{enumerate}\item Pengantar AI, agen cerdas, dan algoritma pencarian\item Representasi pengetahuan, logika, dan sistem inferensi\item Sistem pakar, Jaringan Bayes, dan probabilistic reasoning\item Machine learning, NLP dasar, dan Computer Vision dasar\end{enumerate} &  &  &  \\
\SetCell[r=2]{l} Pustaka & \SetCell[c=4]{l} \textbf{Utama:}\begin{enumerate}\item Russell, S. \& Norvig, P. 2021. Artificial Intelligence: A Modern Approach (4th ed.). Pearson.\item Ertel, W. 2017. Introduction to Artificial Intelligence. Springer.\item Nilsson, N. 2010. The Quest for Artificial Intelligence. Cambridge University Press.\end{enumerate} &  &  &  \\
 & \SetCell[c=4]{l} \textbf{Pendukung:} Materi LMS, dokumentasi resmi, dan jobsheet praktikum. &  &  &  \\
Media Pembelajaran & \SetCell[c=2]{l} \textbf{Software:} Python, Jupyter Notebook, LMS. &  & \SetCell[c=2]{l} \textbf{Hardware:} Komputer/laptop, proyektor. &  \\
Nama Dosen Pengampu & \SetCell[c=4]{l} Tim Pengajar Program Studi D4 TI &  &  &  \\
Matakuliah Syarat & \SetCell[c=4]{l} - &  &  &  \\
\end{longtblr}
\begin{longtblr}{width=\textwidth,colspec={|Q[c,m,0.06\textwidth]|X[1.35,l,m]|X[1.08,l,m]|X[1.16,l,m]|X[0.7,l,m]|X[1.24,l,m]|X[1.06,l,m]|X[0.98,l,m]|Q[c,m,0.05\textwidth]|},hlines,vlines,hspan=minimal,rowhead=2,row{1,2}={bg=headgreen,font=\bfseries},rowsep=1pt,colsep=1.5pt}
Mgg.\par Ke & Kemampuan Akhir Yang Direncanakan (Sub-CP-MK) & Bahan kajian (Materi Pembelajaran) & Bentuk dan Metode Pembelajaran & Estimasi Waktu & Pengalaman Belajar Mahasiswa & Kriteria \& Bentuk Penilaian & Indikator Penilaian & Bobot Penilaian (\%) \\
(1) & (2) & (3) & (4) & (5) & (6) & (7) & (8) & (9) \\
1 & SCPMK0706-03101 & Pengantar AI: sejarah, definisi, dan ruang lingkup. & Modalitas: Blended Learning. Bentuk: Luring/Daring. Metode: Case Method, praktik pemrograman, studi kasus, dan peer review. & 1 x 2 x 50' tatap muka; 1 x 2 x 50' tugas/praktik mandiri. & Mahasiswa mengerjakan aktivitas terarah untuk pengantar AI: sejarah, definisi, dan ruang lingkup dan menunjukkan evidence hasil belajar. & Bentuk penilaian: Pengantar AI: Sejarah, Definisi, dan Ruang Lingkup. Mengacu rubrik RTM. & Ketepatan konsep; kualitas implementasi; validasi dan dokumentasi. & 2 \\
2 & SCPMK0706-03101 & Agen cerdas dan lingkungan. & Modalitas: Blended Learning. Bentuk: Luring/Daring. Metode: Case Method, praktik pemrograman, studi kasus, dan peer review. & 1 x 2 x 50' tatap muka; 1 x 2 x 50' tugas/praktik mandiri. & Mahasiswa mengerjakan aktivitas terarah untuk agen cerdas dan lingkungan dan menunjukkan evidence hasil belajar. & Bentuk penilaian: Agen Cerdas dan Lingkungan. Mengacu rubrik RTM. & Ketepatan konsep; kualitas implementasi; validasi dan dokumentasi. & 2 \\
3 & SCPMK0706-03101 & Algoritma pencarian tak terinformasi: BFS, DFS. & Modalitas: Blended Learning. Bentuk: Luring/Daring. Metode: Case Method, praktik pemrograman, studi kasus, dan peer review. & 1 x 2 x 50' tatap muka; 1 x 2 x 50' tugas/praktik mandiri. & Mahasiswa mengerjakan aktivitas terarah untuk algoritma pencarian tak terinformasi: BFS, DFS dan menunjukkan evidence hasil belajar. & Bentuk penilaian: Algoritma Pencarian Tak Terinformasi: BFS, DFS. Mengacu rubrik RTM. & Ketepatan konsep; kualitas implementasi; validasi dan dokumentasi. & 3 \\
4 & SCPMK0706-03101 & Algoritma pencarian terinformasi: A*, heuristik. & Modalitas: Blended Learning. Bentuk: Luring/Daring. Metode: Case Method, praktik pemrograman, studi kasus, dan peer review. & 1 x 2 x 50' tatap muka; 1 x 2 x 50' tugas/praktik mandiri. & Mahasiswa mengerjakan aktivitas terarah untuk algoritma pencarian terinformasi: A*, heuristik dan menunjukkan evidence hasil belajar. & Bentuk penilaian: Algoritma Pencarian Terinformasi: A*, Heuristik. Mengacu rubrik RTM. & Ketepatan konsep; kualitas implementasi; validasi dan dokumentasi. & 3 \\
5 & SCPMK0706-03101 & Kuis 1: pencarian dan agen. & Modalitas: Blended Learning. Bentuk: Luring/Daring. Metode: Case Method, praktik pemrograman, studi kasus, dan peer review. & 1 x 2 x 50' tatap muka; 1 x 2 x 50' tugas/praktik mandiri. & Mahasiswa mengerjakan aktivitas terarah untuk kuis 1: pencarian dan agen dan menunjukkan evidence hasil belajar. & Bentuk penilaian: Kuis 1: Pencarian dan Agen. Mengacu rubrik RTM. & Ketepatan konsep; kualitas implementasi; validasi dan dokumentasi. & 5 \\
6 & SCPMK0706-03101 & Representasi pengetahuan: logika proposisional. & Modalitas: Blended Learning. Bentuk: Luring/Daring. Metode: Case Method, praktik pemrograman, studi kasus, dan peer review. & 1 x 2 x 50' tatap muka; 1 x 2 x 50' tugas/praktik mandiri. & Mahasiswa mengerjakan aktivitas terarah untuk representasi pengetahuan: logika proposisional dan menunjukkan evidence hasil belajar. & Bentuk penilaian: Representasi Pengetahuan: Logika Proposisional. Mengacu rubrik RTM. & Ketepatan konsep; kualitas implementasi; validasi dan dokumentasi. & 3 \\
7 & SCPMK0706-03101 & Representasi pengetahuan: logika predikat. & Modalitas: Blended Learning. Bentuk: Luring/Daring. Metode: Case Method, praktik pemrograman, studi kasus, dan peer review. & 1 x 2 x 50' tatap muka; 1 x 2 x 50' tugas/praktik mandiri. & Mahasiswa mengerjakan aktivitas terarah untuk representasi pengetahuan: logika predikat dan menunjukkan evidence hasil belajar. & Bentuk penilaian: Representasi Pengetahuan: Logika Predikat. Mengacu rubrik RTM. & Ketepatan konsep; kualitas implementasi; validasi dan dokumentasi. & 3 \\
8 & SCPMK0706-03101, SCPMK1005-03102 & UTS: pencarian dan representasi pengetahuan. & Modalitas: Blended Learning. Bentuk: Luring/Daring. Metode: Case Method, praktik pemrograman, studi kasus, dan peer review. & 1 x 2 x 50' tatap muka; 1 x 2 x 50' tugas/praktik mandiri. & Mahasiswa mengerjakan aktivitas terarah untuk UTS: pencarian dan representasi pengetahuan dan menunjukkan evidence hasil belajar. & Bentuk penilaian: UTS: Pencarian dan Representasi Pengetahuan. Mengacu rubrik RTM. & Ketepatan konsep; kualitas implementasi; validasi dan dokumentasi. & 15 \\
9 & SCPMK0706-03101 & Sistem pakar dan mesin inferensi berbasis aturan. & Modalitas: Blended Learning. Bentuk: Luring/Daring. Metode: Case Method, praktik pemrograman, studi kasus, dan peer review. & 1 x 2 x 50' tatap muka; 1 x 2 x 50' tugas/praktik mandiri. & Mahasiswa mengerjakan aktivitas terarah untuk sistem pakar dan mesin inferensi berbasis aturan dan menunjukkan evidence hasil belajar. & Bentuk penilaian: Sistem Pakar dan Mesin Inferensi Berbasis Aturan. Mengacu rubrik RTM. & Ketepatan konsep; kualitas implementasi; validasi dan dokumentasi. & 3 \\
10 & SCPMK1005-03102 & Jaringan Bayes dan probabilistic reasoning. & Modalitas: Blended Learning. Bentuk: Luring/Daring. Metode: Case Method, praktik pemrograman, studi kasus, dan peer review. & 1 x 2 x 50' tatap muka; 1 x 2 x 50' tugas/praktik mandiri. & Mahasiswa mengerjakan aktivitas terarah untuk jaringan Bayes dan probabilistic reasoning dan menunjukkan evidence hasil belajar. & Bentuk penilaian: Jaringan Bayes dan Probabilistic Reasoning. Mengacu rubrik RTM. & Ketepatan konsep; kualitas implementasi; validasi dan dokumentasi. & 3 \\
11 & SCPMK0706-03101 & Pengantar ML: supervised learning. & Modalitas: Blended Learning. Bentuk: Luring/Daring. Metode: Case Method, praktik pemrograman, studi kasus, dan peer review. & 1 x 2 x 50' tatap muka; 1 x 2 x 50' tugas/praktik mandiri. & Mahasiswa mengerjakan aktivitas terarah untuk pengantar ML: supervised learning dan menunjukkan evidence hasil belajar. & Bentuk penilaian: Pengantar ML: Supervised Learning. Mengacu rubrik RTM. & Ketepatan konsep; kualitas implementasi; validasi dan dokumentasi. & 3 \\
12 & SCPMK1005-03102 & Evaluasi model ML: metrik dan validasi. & Modalitas: Blended Learning. Bentuk: Luring/Daring. Metode: Case Method, praktik pemrograman, studi kasus, dan peer review. & 1 x 2 x 50' tatap muka; 1 x 2 x 50' tugas/praktik mandiri. & Mahasiswa mengerjakan aktivitas terarah untuk evaluasi model ML: metrik dan validasi dan menunjukkan evidence hasil belajar. & Bentuk penilaian: Evaluasi Model ML: Metrik dan Validasi. Mengacu rubrik RTM. & Ketepatan konsep; kualitas implementasi; validasi dan dokumentasi. & 3 \\
13 & SCPMK1005-03102 & Kuis 2: sistem pakar dan ML. & Modalitas: Blended Learning. Bentuk: Luring/Daring. Metode: Case Method, praktik pemrograman, studi kasus, dan peer review. & 1 x 2 x 50' tatap muka; 1 x 2 x 50' tugas/praktik mandiri. & Mahasiswa mengerjakan aktivitas terarah untuk kuis 2: sistem pakar dan ML dan menunjukkan evidence hasil belajar. & Bentuk penilaian: Kuis 2: Sistem Pakar dan ML. Mengacu rubrik RTM. & Ketepatan konsep; kualitas implementasi; validasi dan dokumentasi. & 5 \\
14 & SCPMK0706-03101 & Natural Language Processing dasar. & Modalitas: Blended Learning. Bentuk: Luring/Daring. Metode: Case Method, praktik pemrograman, studi kasus, dan peer review. & 1 x 2 x 50' tatap muka; 1 x 2 x 50' tugas/praktik mandiri. & Mahasiswa mengerjakan aktivitas terarah untuk Natural Language Processing dasar dan menunjukkan evidence hasil belajar. & Bentuk penilaian: Natural Language Processing Dasar. Mengacu rubrik RTM. & Ketepatan konsep; kualitas implementasi; validasi dan dokumentasi. & 3 \\
15 & SCPMK0706-03101 & Computer Vision dasar. & Modalitas: Blended Learning. Bentuk: Luring/Daring. Metode: Case Method, praktik pemrograman, studi kasus, dan peer review. & 1 x 2 x 50' tatap muka; 1 x 2 x 50' tugas/praktik mandiri. & Mahasiswa mengerjakan aktivitas terarah untuk Computer Vision dasar dan menunjukkan evidence hasil belajar. & Bentuk penilaian: Computer Vision Dasar. Mengacu rubrik RTM. & Ketepatan konsep; kualitas implementasi; validasi dan dokumentasi. & 3 \\
16 & SCPMK0706-03101, SCPMK1005-03102 & PBL: implementasi sistem AI. & Modalitas: Blended Learning. Bentuk: Luring/Daring. Metode: Case Method, praktik pemrograman, studi kasus, dan peer review. & 1 x 2 x 50' tatap muka; 1 x 2 x 50' tugas/praktik mandiri. & Mahasiswa mengerjakan aktivitas terarah untuk PBL: implementasi sistem AI dan menunjukkan evidence hasil belajar. & Bentuk penilaian: PBL: Implementasi Sistem AI. Mengacu rubrik RTM. & Ketepatan konsep; kualitas implementasi; validasi dan dokumentasi. & 22 \\
17 & SCPMK1005-03102 & UAS: presentasi dan evaluasi sistem AI. & Modalitas: Blended Learning. Bentuk: Luring/Daring. Metode: Case Method, praktik pemrograman, studi kasus, dan peer review. & 1 x 2 x 50' tatap muka; 1 x 2 x 50' tugas/praktik mandiri. & Mahasiswa mengerjakan aktivitas terarah untuk UAS: presentasi dan evaluasi sistem AI dan menunjukkan evidence hasil belajar. & Bentuk penilaian: UAS: Presentasi dan Evaluasi Sistem AI. Mengacu rubrik RTM. & Ketepatan konsep; kualitas implementasi; validasi dan dokumentasi. & 23 \\
\end{longtblr}
